\begin{textsample}{NEG Dim 1 – human – Score -6.00 – rj/rj\_ef\_linguagens\_ed\_fis\_2\_p10.txt}  \label{ex:f1_neg_009}
A organização das Unidades Temáticas se baseia na compreensão de que todas as práticas \textbf{corporais} são permeadas pelo caráter lúdico, embora essa não seja a finalidade da Educação Física ( BNCC, 2017, p. 218 ).
Para além da ludicidade, ao praticar qualquer atividade \textbf{corporal} proposta pelas Unidades Temáticas, os estudantes podem se apropriar de lógicas intrínsecas como regras, códigos e táticas, o que vai se relacionar com representações e \textbf{significados} atribuídos às relações sociais de \textbf{maneira} geral.
Por essa razão, a delimitação das habilidades privilegia oito dimensões de conhecimento ( Experimentação, Uso e Apropriação, Fruição, Reflexão sobre a Ação, Construção de Valores, Análise, Compreensão e \textbf{Protagonismo} Comunitário ), relacionando -se com as Competências Gerais da Educação Básica e Específicas da área da Educação Física no Ensino Fundamental.

% matched lemmas: corporal, maneira, protagonismo, significado
\end{textsample}
